\section{目的}

本研究の目的は、Webブラウザ上のエッジ環境においてリアルタイムかつ安定して動作する待機列のグループ検出手法を開発し、その検出精度とシステム性能を段階的に向上・評価することである。
具体的には、前節で挙げた課題を解決するため、以下の4つの小目的を設定し検証を行う。

\begin{enumerate}
    \item \textbf{多角的な特徴量（ポーズ・体の向き・幾何関係）を活用したグループ判定精度の向上}：
    単一の二次元バウンディングボックスによる評価を超え、顔・骨格情報から得られる体の向きや空間的距離、前後の並び関係を総合的に考慮したグループ判定アルゴリズムを開発し、密集による遮蔽環境や多様な待機形態における誤判定を削減する。

    \item \textbf{複数フレーム情報を利用した空間的・時間的トラッキングの実現}：
    時間軸に沿った人物追跡およびカルマンフィルタ等を活用した動き解析を組み込むことで、一時的な通行人やノイズによる影響を排除し、待機列内に滞留する真の構成員のみを高精度に追跡・グループ化する。

    \item \textbf{エッジ環境（ブラウザ）に最適化した堅牢かつ高効率な推論パイプラインの構築}：
    Web Workerによるマルチスレッド処理、キャンバスやバッファの再利用によるメモリ解放の厳格化、および最新の軽量高精度モデル（YOLOv12 Nano等）の導入により、OOM（Out of Memory）を完全に防止しつつ実用的なフレームレートを達成する。

    \item \textbf{全13段階のアプローチに対する比較検証と評価手法の提示}：
    ベースラインから最終的な特徴量整理に至る全13種類のアプローチを個別に実装・評価し、各段階における設計思想、定量・定性成果、計算負荷、および残された課題を明示することで、将来的な待機列制御システム構築における要素技術の選定指針を提示する。
\end{enumerate}